\documentclass[11pt,letterpaper]{article}

\usepackage[T1]{fontenc}
\usepackage[utf8]{inputenc}
\usepackage{lmodern}
\usepackage[margin=0.82in]{geometry}
\usepackage{amsmath,amssymb,amsthm,amscd,mathtools}
\usepackage{booktabs,longtable,array}
\usepackage{enumitem}
\usepackage{xcolor}
\usepackage{microtype}
\usepackage[hidelinks]{hyperref}
\usepackage[nameinlink,noabbrev]{cleveref}

\definecolor{Verified}{HTML}{176B3A}
\definecolor{Conditional}{HTML}{9A6700}
\definecolor{Open}{HTML}{B42318}
\definecolor{Refuted}{HTML}{7A271A}
\definecolor{Provenance}{HTML}{475467}

\newcommand{\Status}[2]{\mbox{\textcolor{#1}{\small\textbf{[#2]}}}}
\newcommand{\Verified}{\Status{Verified}{VERIFIED}}
\newcommand{\Conditional}{\Status{Conditional}{CONDITIONAL}}
\newcommand{\Open}{\Status{Open}{OPEN}}
\newcommand{\Refuted}{\Status{Refuted}{REFUTED}}
\newcommand{\Source}[1]{\textcolor{Provenance}{\nolinkurl{#1}}}

\newtheorem{theorem}{Theorem}[section]
\newtheorem{proposition}[theorem]{Proposition}
\newtheorem{lemma}[theorem]{Lemma}
\newtheorem{corollary}[theorem]{Corollary}
\theoremstyle{definition}
\newtheorem{obligation}[theorem]{Proof obligation}
\newtheorem{test}[theorem]{Disproof test}
\theoremstyle{remark}
\newtheorem{remark}[theorem]{Remark}
% Explicit names keep cleveref well-defined for the custom environments.
\crefname{obligation}{proof obligation}{proof obligations}
\Crefname{obligation}{Proof obligation}{Proof obligations}
\crefname{test}{disproof test}{disproof tests}
\Crefname{test}{Disproof test}{Disproof tests}

\title{S6 Proof Repair and Disproof Workbench --- Private}
\author{Working mathematical notebook following the three-stream proof audit}
\date{26 August 2026}

\begin{document}
\maketitle

\begin{abstract}
This is a separate mathematical workbench for validating, repairing, or
disproving the reconstructed compact complex threefold on $S^6$.  The finite
integer calculations are recorded as exact certificates: the monodromy
invariant and coinvariant lattices, the finite-fibre indices, the parameter
$p=12\ell_0-4\ell_1-3\ell_2$ and Smith form
$\operatorname{diag}(1,|p|)$, and the degree-two matrix with primitive
left kernel $\Phi=(4,2,3,2)$.  The global conclusions remain
\Conditional{} until the printed matrices are identified with the actual
geometric specialization, attaching, Mayer--Vietoris, and Leray maps.

The notebook names those bridge lemmas (C1--C3, V1, MV1, and L1), gives a
conditional recognition theorem, and treats the Section~10/CDP issue in
three separate questions: is a printed proof step valid, is the published
theorem false, and does this candidate satisfy its hypotheses?  It does not
silently alter the reconstructed manuscript; the frozen 121-page audit and
108-page transcription remain under \Source{audit_baseline/}.

For the Section~10 route, the source proof of the candidate's compact smooth
analytic geometry and full cusp fibre has now been independently retraced.
The conductor-supported differential, compact-Cartier-fibre duality, and
top-degree base change then prove
$R^2f_*(T_X\otimes L)\ne0$ for every
$L\in\operatorname{Pic}^0(X)$, without assuming the unresolved equality
$H^2(X;\mathbb Z)=0$.  Extension to every component of
$\operatorname{Pic}(X)$ is equivalent for this method to the separate open
restriction-map statement $\eta^*i^*=0$.  None of this promotes the global
topology or by itself declares the published CDP theorem false.
\end{abstract}

\tableofcontents

\section{Problem, evidence boundary, and status language}

The target is the manuscript's claimed compact complex threefold $X$ equipped with a holomorphic map $f\colon X\to\mathbb P^1$, with
\[
X\cong_{\mathrm{diff}}S^6,\qquad a(X)=1,\qquad
\langle c_3(T_X),[X]\rangle=2.
\]
The review found no finite-algebra contradiction in the completed manuscript, but it did not certify the geometric inputs from which the terminal matrices are obtained.

\subsection{Status rules}

\begin{description}[leftmargin=8.4em,style=nextline]
\item[\Verified] Reproduced exactly or proved in this notebook from stated hypotheses.
\item[\Conditional] Correct after explicitly named geometric or analytic hypotheses are granted.
\item[\Open] A load-bearing statement for which no complete proof or counterexample is yet recorded here.
\item[\Refuted] A statement defeated by a printed counterexample or a complete contradiction.
\end{description}

\subsection{Current ledger}

\begin{longtable}{>{\raggedright\arraybackslash}p{0.20\textwidth}>{\raggedright\arraybackslash}p{0.18\textwidth}>{\raggedright\arraybackslash}p{0.53\textwidth}}
\toprule
Status & Object & Present evidence \\
\midrule
\endhead
\Verified & terminal finite algebra & Monodromy products/orders, Smith forms, exterior-power ranks, displayed lattice saturation, period identities, $\eta^2=12$, and the Gram determinant have exact checks in the retained review harnesses. \\
\Verified & candidate analytic geometry & Sections~3--6 construct the smooth compact connected threefold, proper holomorphic map, full reduced normal-crossings cusp fibre, $\mathrm{dP}_6$ normalization, and a trivial canonical bundle on its open cusp neighbourhood.  The source proof chain has been independently retraced. \\
\Verified & conductor differential & The normalization--conductor sequence and the local $x$, $xy$, and $xyz$ models produce the nonzero ambient differential under the stated reduced normal-crossings and pullback-triviality hypotheses. \\
\Verified & compact-Cartier duality & Ramis--Ruget duality, its Gorenstein complete-intersection calculation, Cartier adjunction, and compact coherent-cohomology finiteness give the stated perfect pairing without a normality hypothesis. \\
\Verified & top-degree base change & The analytic cohomology-and-base-change theorem gives $\beta^2_c$ as an isomorphism when $H^3$ of the fibre vanishes; compactness, smoothness, connectedness, and nonconstancy make properness and flatness automatic over a smooth curve. \\
\Verified & Section~10 on $\operatorname{Pic}^0$ & Local canonical triviality and $c_1(L)=0$ show $\eta^*A_L\simeq\mathcal O_{\mathrm{dP}_6}$ for every $L\in\operatorname{Pic}^0(X)$.  Conductor descent, duality, and base change give $R^2f_*(T_X\otimes L)\ne0$ for every such $L$. \\
\Conditional & $\pi_1(X)$ and $H_*(X;\mathbb Z)$ & Correct for the printed presentations and boundary matrices; their equality with the maps induced by the analytic gluings is not yet established here. \\
\Open & cusp collapse & The analytic action, quotient, and central-fibre charts are verified.  C1 still requires a finite equivariant CW model and a global strong deformation retraction compatible with all strata. \\
\Open & integral topology maps & Specialization, transfer indices, saturation, torsion, primitive bases, orientations, and degree-three signs. \\
\Open & all-line-bundle extension & For this conductor method, extension from $\operatorname{Pic}^0(X)$ to all of $\operatorname{Pic}(X)$ is equivalent to $r=\eta^*i^*=0$.  The claimed $H^2(X;\mathbb Z)=0$ is sufficient but remains conditional on the topology bridges. \\
\Open & CDP adjudication & The printed cotangent-sheaf reduction misses the conductor-supported section.  A line-by-line repair search and the candidate's remaining global CDP hypotheses must be settled before declaring Theorem~2.2 false. \\
\bottomrule
\end{longtable}

The working rule is asymmetric: a failed proof step does not refute the theorem, and a coherent terminal calculation does not prove that the required geometric map exists.

\section{Fixed algebraic data reproduced by the audits}

This section freezes the exact statements that need no further numerical experimentation. The research task is to prove that the geometric maps have these algebraic representatives.

Let
\[
H_1(F;\mathbb Z)=\Lambda
=\langle\widehat\gamma,\widehat u,\widehat w,\widehat\delta\rangle,
\qquad
H^1(F;\mathbb Z)=V
=\langle\gamma,u,w,\delta\rangle.
\]

\begin{proposition}[checked monodromy lattices]
\Verified\ The displayed monodromy matrices satisfy
\[
A_1A_2M_0=I,
\qquad
\ker(M_0-I)=\operatorname{im}(M_0-I)
=\Lambda_{\mathrm{tor}}:=\langle\widehat w,\widehat\delta\rangle.
\]
For $T_0=I+N$ with $N^2=0$, the invariant lattices in degrees one through three are
\begin{align*}
V^{T_0}&=\langle\gamma,u\rangle,\\
(\wedge^2V)^{T_0}
&=\langle\gamma u,\gamma\delta,uw,\gamma w-u\delta\rangle,\\
(\wedge^3V)^{T_0}
&=\langle\gamma uw,\gamma u\delta\rangle.
\end{align*}
They are saturated and have ranks $(1,2,4,2,1)$ across all degrees. For $G=\langle T_1,T_2\rangle$, the nonzero coinvariant lattices used downstream are
\[
V_G=\mathbb Z\bar\delta,\qquad
(\wedge^2V)_G=\mathbb Z[\gamma\delta],\qquad
(\wedge^3V)_G=\mathbb Z[uw\delta],
\]
all torsion-free, with $[uw]=6[\gamma\delta]$.
\end{proposition}

\begin{proposition}[checked finite-fibre indices]
\Verified\ In degrees $q=0,1,2,3,4$, the two finite-fibre
specialization-index tuples in the manuscript's bases are
\[
(1,3,1,1,3),\qquad (1,4,2,2,4).
\]
For $v_2=-\varepsilon'$, the degree-three index-two quotient at $p_2$ is generated by $[\gamma uw]$.
\end{proposition}

\begin{proposition}[checked group presentation]
\Verified\ Once the asserted based inclusions are granted, the presentation reduces to
\[
\left\langle c,x,y\ \middle|\ c\ \text{central},\ xy=c^{\ell_0},\
x^3=c^{\ell_1},\ y^4=c^{\ell_2}\right\rangle.
\]
After eliminating $y$, its relation matrix is
\[
R=\begin{pmatrix}3&-\ell_1\\4&\ell_2-4\ell_0\end{pmatrix},
\qquad
\det R=-p,
\qquad
p=12\ell_0-4\ell_1-3\ell_2,
\]
and its Smith form is $\operatorname{diag}(1,|p|)$.
\end{proposition}

\begin{proposition}[checked degree-two terminal matrix]
\Verified\ The displayed degree-two matrix has columns
\[
(2,-4,-2,3),\ (1,-2,-2,3),\ (3,0,-4,0),\
(0,1,0,-1),\ 0,\ 0.
\]
Its image has rank three and is saturated, its cokernel is $\mathbb Z$, and its primitive left-kernel generator is
\[
\Phi=(4,2,3,2).
\]
\end{proposition}

\begin{proposition}[checked all-degree terminal homology package]
\Verified\ For $p\ne0$, the retained exact audit applied to the manuscript's
printed degree-by-degree Mayer--Vietoris/Leray data gives the terminal
abelian-group tuple
\[
 \bigl(H_0,H_1,H_2,H_3,H_4,H_5,H_6\bigr)
 \cong
 \bigl(\mathbb Z,\mathbb Z/|p|\mathbb Z,
 \mathbb Z/|p|\mathbb Z,\mathbb Z/|p|\mathbb Z,
 \mathbb Z/|p|\mathbb Z,0,\mathbb Z\bigr).
\]
Consequently the printed algebraic package has the homology of $S^6$ when
$p=-1$.  The exact evidence is frozen in
\Source{audit_baseline/annotated_audit.tex}, ledger item B-EX-06, and in the
degree-by-degree calculation surrounding Theorem~7.22 of
\Source{audit_baseline/clean_transcription.tex}.
\end{proposition}

Every proposition in this section concerns the printed matrices and lattices. None identifies those matrices with a geometric specialization, attaching, Mayer--Vietoris, or Leray map.

\section{The prior Fable certificate: relevance and boundary}

The separate Erd\H{o}s--Strauss workspace contains a pinned copy of Jon
Poplett's JGPTech package for the polynomial map publicly attributed there to
Fable.  This is relevant as a standard of verification, not as evidence for
any step of the present $S^6$ construction.  The canonical checkout and its
immutable mirror are listed with exact absolute paths in
\Source{comparative_context/README.md}; within either checkout the relevant
file is \Source{JGPTech_Fun/Jacobian/aceg.py}.

Define $F=(F_1,F_2,F_3)\colon\mathbb C^3\to\mathbb C^3$ by
\begin{align*}
F_1(x,y,t)&=(1+xy)^3t+y^2(1+xy)(4+3xy),\\
F_2(x,y,t)&=y+3x(1+xy)^2t+3xy^2(4+3xy),\\
F_3(x,y,t)&=2x-3x^2y-x^3t.
\end{align*}

\begin{proposition}[replayed finite certificate]
\Verified\ Exact rational-polynomial arithmetic gives
\[
\det DF=-2.
\]
Moreover, the three distinct points
\[
\left(0,0,-\tfrac14\right),\qquad
\left(1,-\tfrac32,\tfrac{13}{2}\right),\qquad
\left(-1,\tfrac32,\tfrac{13}{2}\right)
\]
all map to $(-\tfrac14,0,0)$.  The local pinned checker
\Source{Jacobian/aceg.py} was replayed successfully on 26 August 2026 with
the `base` and manifest checks.  Its exact SHA-256, mirror commit, and
absolute paths are recorded in \Source{comparative_context/README.md}.
Thus the displayed map itself has the advertised constant nonzero Jacobian
and noninjective fibre.
\end{proposition}

\begin{remark}[provenance boundary]
The located corpus is an independently organized proof and exact audit of the
displayed map, not an original model-session transcript.  No original
Fable/Claude model-session source was located.  None of this certificate
supplies a lemma about the complex-threefold manuscript, and the Jacobian
counterexample is kept distinct from Erd\H{o}s--Strauss arithmetic.
\end{remark}

The transferable lesson is procedural.  The prior result terminates in a
small certificate---one determinant identity and one explicit collision---that
can be checked independently of its discovery narrative.  The present paper
will reach the same evidentiary level only when its analytic gluings terminate
in independently reconstructed integral chain maps.  Accordingly this
workbench uses typed claims, logical dependency edges, typed morphisms,
explicit uncertainties, clean-archive replays, and rendered-page QA; see
\Source{comparative_context/README.md}.

\section{Topology: exact reduction and geometric obligations}

\subsection{What the exact arithmetic already proves}

Let
\[
p=12\ell_0-4\ell_1-3\ell_2.
\]
The retained exact harness verifies the determinant and Smith-normal-form calculations for the printed relation matrix. In particular, the manuscript's chosen parameters $(\ell_0,\ell_1,\ell_2)=(0,1,-1)$ give $p=-1$.

\begin{proposition}[terminal group calculation]\label[proposition]{prop:terminal-group}
\Verified\ For the printed based presentation and its stated meridian relations, the resulting group is cyclic of order $|p|$. Hence the printed presentation is trivial for $(0,1,-1)$.
\end{proposition}

\begin{proof}
The exact row and column operations recorded in the audit reduce the displayed
$2\times2$ relation matrix to Smith normal form
$\operatorname{diag}(1,|p|)$. A group presented by that matrix is therefore
$\mathbb Z/|p|\mathbb Z$, which is trivial when $p=-1$.
\end{proof}

\begin{remark}
Proposition~\ref{prop:terminal-group} is a theorem about the printed
presentation. It becomes a theorem about $\pi_1(X)$ only after the based
inclusions induced by all three analytic attaching maps have been derived in
the same generators and orientations.
\end{remark}

\subsection{The map that must be identified}

Write the relevant decomposition schematically as $X=U\cup N$, where $N$ is the union of the filling pieces and $U\cap N$ is the disjoint union of their collars. The homological boundary map is
\[
\delta_q=(i_{U*},-i_{N*})\colon
H_q(U\cap N;\mathbb Z)\longrightarrow
H_q(U;\mathbb Z)\oplus H_q(N;\mathbb Z).
\]
The printed matrices denoted $\alpha_q$ have verified ranks, Smith forms, and saturation properties. The central question is not another Smith-form computation; it is the equality
\[
\boxed{\alpha_q=[\delta_q]_{\mathcal B_q}}
\]
for explicitly constructed primitive oriented integral bases $\mathcal B_q$.

The equality is decomposed below into the six nonredundant proof obligations
\hyperref[obl:C1]{C1}, \hyperref[obl:C2]{C2}, \hyperref[obl:C3]{C3},
\hyperref[obl:V1]{V1}, \hyperref[obl:MV1]{MV1}, and
\hyperref[obl:L1]{L1}.  The first three construct the finite equivariant
model, compute its integral collapse map, and identify that map with
nearby-cycle specialization.  Obligation V1 fixes the based van Kampen
inclusions; MV1 identifies the actual Mayer--Vietoris boundary cylinders and
multiple-fibre indices; L1 supplies an independently oriented
Leray--Mayer--Vietoris comparison.  This crosswalk is exhaustive for the
conditional endpoint stated here; none of the six is replaced by matching a
rank or by choosing the desired matrix as a normalization.

\begin{test}[fast falsification targets]
Any of the following disproves the claimed topology calculation: a nonprimitive geometric image where $\alpha_q$ assumes a primitive one; an index or torsion class omitted by specialization; a meridian sign changing $p=-1$ to a nonunit; or a Leray transgression differing from $\pm p\omega$ by a nonunit factor.
\end{test}

\subsection{Conditional recognition}

\begin{theorem}[precise conditional endpoint]
\Conditional\ Assume obligations \hyperref[obl:C1]{C1}--\hyperref[obl:C3]{C3},
\hyperref[obl:V1]{V1}, \hyperref[obl:MV1]{MV1}, and
\hyperref[obl:L1]{L1} are discharged and give exactly the printed integral
maps. Then the verified terminal calculations yield
\[
\pi_1(X)=0,\qquad H_k(X;\mathbb Z)=
\begin{cases}
\mathbb Z,&k=0,6,\\
0,&1\le k\le5.
\end{cases}
\]
Consequently $X$ is a smooth homotopy $6$-sphere; the standard fact $\Theta_6=0$ identifies it diffeomorphically with $S^6$.
\end{theorem}

\begin{proof}
The hypotheses identify the actual van Kampen and Mayer--Vietoris/Leray maps with the already verified presentations and matrices. Their Smith forms give the displayed fundamental group and homology. A simply connected integral homology sphere is a homotopy sphere by Hurewicz and Whitehead; the final smooth recognition uses $\Theta_6=0$.
\end{proof}

\section{First proof attacks: exact lemmas to establish or break}

The following statements turn the review's broad objections into falsifiable mathematical targets. A proof must derive the maps from the analytic charts; matching ranks is insufficient.

\subsection{Cusp quotient and collapse}

The source defines an infinite toric threefold $Y$ from the height-one cones over the periodic $A_2$ triangulation and an action of the form
\[
\Psi_{\bar\lambda}(x,t)
=\bigl(c_{\bar\lambda}(t)t^{B_0\bar\lambda}x,t\bigr),
\qquad \bar\lambda\in\bar\Lambda.
\]
Theorem~4.5's analytic part has now been independently retraced: this action
is free and properly discontinuous, its quotient is a smooth Hausdorff
threefold proper over the disc, and its central fibre has the asserted local
normal-crossings equations.  C1 below is narrower: it asks for the finite
equivariant topological model and collapse that are not supplied by those
analytic quotient facts.

\begin{obligation}[C1: finite equivariant cusp model]\label[obligation]{obl:C1}
Construct a locally finite $\bar\Lambda$-equivariant CW model of $(Y_{\le\eta},Y_0)$ whose quotient is finite, and prove that the collapse is equivariant in the diagram
\[
\begin{CD}
Y_{\le\eta} @>{r_Y}>> Y_0\\
@VV{/\bar\Lambda}V @VV{/\bar\Lambda}V\\
N'_0 @>{r}>> W.
\end{CD}
\]
The proof may cite the verified analytic action theorem, but it must still
establish continuity of the proposed collapse across the one-, two-, and
three-branch strata, compatibility with the fibre-torus action, finiteness of
the quotient CW model, and a strong deformation retraction.
\end{obligation}

\begin{obligation}[C2: integral collapse matrix]\label[obligation]{obl:C2}
With fixed cellular orientations, prove for $0\le q\le4$ the exact sequence
\[
0\longrightarrow
\operatorname{im}(\wedge^qM_0-I)
\longrightarrow\wedge^q\Lambda
\xrightarrow{c_*}H_q(W;\mathbb Z)
\longrightarrow0.
\]
In particular, derive $c_*\colon\Lambda\to H_1(W)=\bar\Lambda$ as the quotient map, prove
\[
c_*(\widehat\gamma\wedge\widehat u)=[\mathbb T_B],
\qquad
c_*(\widehat w\wedge\widehat\delta)=0,
\]
and record the degree-three orientations rather than inferring them from ranks.
\end{obligation}

\begin{obligation}[C3: collapse equals specialization]\label[obligation]{obl:C3}
Prove functorially and integrally that
\[
c^*=i_F^*r^*=\operatorname{sp}_{\mathrm{nc}}^q,
\qquad
\operatorname{im}c^*=(\wedge^qV)^{T_0},
\]
including saturation, torsion, bases, and the degree-three sign. The comparison diagram must state the precise nearby-cycle edge map; the equality may not be defined by choosing the manuscript's answer as a normalization.
\end{obligation}

\subsection{Based inclusions}

Choose one basepoint on the zero section and fixed whiskers through every collar.

\begin{obligation}[V1: primitive meridians and drift]\label[obligation]{obl:V1}
Derive from the coordinate gluings, with no unrecorded conjugations,
\[
\widehat g_j^{m_j}=t_{v_j},
\qquad
\widetilde a_0=t_\mu\widehat\rho_2^{-1}\widehat\rho_1^{-1}.
\]
Prove that the clockwise meridian lifts to $\widehat g_j$ rather than its inverse or a nonprimitive power, that the two finite-fibre signs agree with the chosen orientations, that the zero-section whiskers add no hidden translation, and that the tautological cusp disc gives $\mu=0$.
\end{obligation}

\subsection{Mayer--Vietoris and Leray}

\begin{obligation}[MV1: geometric boundary cylinders]\label[obligation]{obl:MV1}
Using the independently constructed collars and the maps from C2--C3, prove
for every relevant degree $q$ that
\[
 \alpha_q=
 \bigl[(i_{U*},-i_{N*}):H_q(U\cap N;\mathbb Z)
 \longrightarrow H_q(U;\mathbb Z)\oplus H_q(N;\mathbb Z)\bigr]_{\mathcal B_q}
\]
in the manuscript's named primitive oriented bases.  This must include a
chain-level check of every zero column, sign, torsion class, and saturation
claim; equality of ranks or Smith forms is not sufficient.  For
$K_k=\ker(H_k(M)\to H_k(W))$, derive geometrically
\[
K_1=\mathbb Z[\widetilde a_0],
\qquad
K_2=\langle\xi_{\widehat w},\xi_{\widehat\delta}\rangle,
\]
and fix one orientation convention for
\[
\partial(j_*\xi_\nu)=(I-A_1^{-1})\nu=(I-A_2)\nu.
\]
The same calculation must recover the multiple-fibre class
\([F]\longmapsto(3[S_1],-4[S_2])\) and the index-two degree-three image at $p_2$.
\end{obligation}

\begin{obligation}[L1: transgression comparison]\label[obligation]{obl:L1}
Construct the clutching cocycle that compares the Leray transgression with the Mayer--Vietoris obstruction. For $12\gamma$, derive the discrepancy triple
\[
\bigl(4\ell_1-\alpha_1,\ 3\ell_2-\alpha_2,\
-(12\ell_0-\alpha_1-\alpha_2)\bigr),
\]
whose sum is $-p$, and thereby determine---rather than assume---the sign in
\[
d_2^{0,1}(12\gamma)=\pm p\,\omega.
\]
Then prove that the cup-product and coinvariant identifications are natural in the same orientations before using multiplicativity in higher degrees.
\end{obligation}

Until C1--C3, V1, MV1, and L1 are established, the conclusion $X\cong S^6$ remains conditional. The integer algebra downstream of these statements is already fixed by the previous section.

\section{The non-normal fibre: exact claim under investigation}

\paragraph{Status convention.}
This notebook does not assert that the CDP theorem is false. It studies one
specific inference in the printed proof and the conditional application of a
conductor-supported differential to the candidate fibre $W$. A defect in the
printed inference, an alternative proof of the published theorem, and failure
of the candidate construction are logically distinct possibilities.

\subsection{The local conductor proposition}

\begin{proposition}[conductor-supported ambient differential]
\label[proposition]{prop:conductor-ambient}
\Verified\ Let $g:\mathcal X\to\Delta$ be proper, with $\mathcal X$ a
smooth complex threefold, and let $S=g^{-1}(0)$ be a reduced
normal-crossings surface. Write $\eta:\widetilde S\to S$ for the
normalization, $D\subset S$ for the reduced singular locus, and
$\widetilde D=\eta^{-1}(D)$. Let $A$ be a line bundle on $S$ such that
$\eta^*A\simeq\mathcal O_{\widetilde S}$, and choose a nowhere-vanishing
section $e$ of $\eta^*A$. Then there is a nonzero section
\[
 s\in H^0\!\left(S,\Omega^1_{\mathcal X}|_S\otimes A\right)
\]
whose image in $\Omega^1_S\otimes A$ is supported on $D$, hence is
torsion, and whose image in
\[
 \widetilde\Omega^1_S\otimes A,
 \qquad
 \widetilde\Omega^1_S:=\Omega^1_S/\mathcal T(\Omega^1_S),
\]
is zero. If $S$ is connected and $A\not\simeq\mathcal O_S$, then
$H^0(S,N^*_{S/\mathcal X}\otimes A)=H^0(S,A)=0$, so $s$ does not come
from the conormal term.
\end{proposition}

\begin{proof}
Locally $g=ux$, $uxy$, or $uxyz$, where $u$ is a unit. The conormal map is
generated by $dg|_S$. At a double point,
\[
 dg|_S=u(y\,dx+x\,dy),
\]
and at a triple point it has coefficients in $(xy,xz,yz)$. Thus $dg|_S$
is nonzero away from $D$ but restricts to zero on $D$; equivalently, it
lies in the conductor ideal times $\Omega^1_{\mathcal X}|_S$.

The normalization--conductor square gives
\[
0\longrightarrow\mathcal O_S\longrightarrow
 \eta_*\mathcal O_{\widetilde S}\oplus\mathcal O_D
 \longrightarrow\eta_*\mathcal O_{\widetilde D}\longrightarrow0.
\]
After tensoring with the locally free sheaf
$E=\Omega^1_{\mathcal X}|_S\otimes A$, the compatible pair
\[
 \bigl(\eta^*(dg|_S)\otimes e,\,0\bigr)
\]
descends uniquely to $s$. It is nonzero off $D$. Its intrinsic differential
vanishes off $D$, so its image in $\Omega^1_S\otimes A$ is
conductor-supported torsion. For nontrivial $A$, a global section of $A$
pulls back to constants on the compact connected components of
$\widetilde S$. Compatibility across the conductor and connectedness of
$S$ force either all constants to vanish or all to be nonzero; the latter
would trivialize $A$. Hence $H^0(S,A)=0$.
\end{proof}

Consequently the implication
\[
\left[
 H^0(S,N^*_{S/\mathcal X}\otimes A)=0,\quad
 H^0(S,\widetilde\Omega^1_S\otimes A)=0
\right]
\Longrightarrow
H^0(S,\Omega^1_{\mathcal X}|_S\otimes A)=0
\]
is not a formal consequence of the conormal sequence for a non-normal
normal-crossings surface. The official CDP article prints precisely this
reduction on p.~692 \cite[p.~692]{CDP2020}. This observation concerns that inference, not by itself
the truth value of the theorem.

\begin{lemma}[pullback triviality on a $\mathrm{dP}_6$ normalization]
\label[lemma]{lem:dp6-pullback}
\Verified\ Let $i:W\hookrightarrow X$ be a closed complex subspace with
$H^2(X;\mathbb Z)=0$, and suppose the normalization
$\eta:\widetilde W\to W$ satisfies
$\widetilde W\simeq\mathrm{dP}_6$.  Then, for every
$B\in\operatorname{Pic}(X)$,
\[
 \eta^*i^*B\simeq\mathcal O_{\widetilde W}.
\]
\end{lemma}

\begin{proof}
The group containing $c_1(B)$ is zero, so
$c_1(\eta^*i^*B)=0$.  Since
$H^1(\mathrm{dP}_6,\mathcal O)=0$, the exponential sequence makes
\[
 c_1:\operatorname{Pic}(\mathrm{dP}_6)
 \longrightarrow H^2(\mathrm{dP}_6;\mathbb Z)
\]
injective.  Hence the pullback line bundle is trivial.
\end{proof}

\subsection{Candidate geometry and the topology-free
\texorpdfstring{$\operatorname{Pic}^0$}{Pic0} route}

\begin{proposition}[analytic geometry of the reconstructed cusp fibre]
\label[proposition]{prop:candidate-geometry}
\Verified\ The reconstructed source produces a connected compact smooth
complex threefold $X$, a proper surjective holomorphic map
$f:X\to B\simeq\mathbb P^1$, and a full scheme-theoretic fibre
$W=f^{-1}(p_0)$ with the following properties:
\begin{enumerate}[label=(\roman*)]
\item $W$ is a connected reduced normal-crossings Cartier surface, with
local equations $z_0$, $z_0z_1$, and $z_0z_1z_2$ at its one-, two-, and
three-branch points;
\item its normalization is $\eta:\mathrm{dP}_6\to W$;
\item $W$ has an open neighbourhood $N_0\subset X$ on which
$\omega_X$ is trivial.
\end{enumerate}
\end{proposition}

\begin{proof}
The period construction in Theorem~3.4 of the frozen transcription supplies
the admissible family and the cusp-holomorphic matrix $C(t)$.  In source
pp.~26--28 (lines 1331--1446), Theorem~4.5 proves that the toric action is
free and properly discontinuous, that its quotient $N_0$ is a connected
Hausdorff smooth threefold, that $f_0:N_0\to\Delta$ is proper, and that its
full fibre is reduced with exactly the three displayed local equations.  The
same proof constructs an invariant nowhere-vanishing holomorphic three-form
on $N_0$.

Proposition~4.6, source p.~29 (lines 1453--1465), proves that the finite map
$\mathrm{dP}_6\to W$ is surjective and generically injective and hence is
the normalization.  Proposition~4.7(i)--(iii) and Theorem~5.4 give the
holomorphic overlap maps.  Finally Theorem~6.2, source pp.~36--37
(lines 1811--1834), checks the quotient atlas, Hausdorffness, properness,
compactness, and connectedness, and identifies
$f^{-1}(D_0)=N_0$ and $f^{-1}(p_0)=W$ as a reduced divisor.  Because the
embedding $N_0\hookrightarrow X$ is open and holomorphic, the constructed
trivialization of $K_{N_0}$ is a trivialization of $\omega_X|_{N_0}$.

These source gates were rechecked independently in the present workbench;
no topology, Mayer--Vietoris, or Leray assertion is used here.
\end{proof}

\begin{proposition}[$\operatorname{Pic}^0$ pullback without global topology]
\label[proposition]{prop:pic0-pullback}
\Verified\ Let $X$ be a compact complex manifold, let
$i:W\hookrightarrow X$ have normalization
$\eta:\widetilde W\to W$ with $\widetilde W\simeq\mathrm{dP}_6$, and
suppose $\omega_X$ is trivial on a neighbourhood of $W$.  For every
$L\in\operatorname{Pic}^0(X)$, put
\[
 A_L=(L^\vee\otimes\omega_X)|_W.
\]
Then
\[
 \eta^*A_L\simeq\mathcal O_{\widetilde W}.
\]
\end{proposition}

\begin{proof}
Every $L\in\operatorname{Pic}^0(X)$ has $c_1(L)=0$.  Therefore
$c_1(\eta^*i^*L)=0$.  Since
$H^1(\mathrm{dP}_6,\mathcal O)=0$, the exponential sequence makes the
Chern-class map on $\operatorname{Pic}(\mathrm{dP}_6)$ injective, so
$\eta^*i^*L$ is trivial.  The neighbourhood trivialization of $\omega_X$
also makes $\eta^*i^*\omega_X$ trivial, proving the formula for $A_L$.
\end{proof}

For the stronger all-line-bundle assertion, define
\[
 r:=\eta^*i^*:\operatorname{Pic}(X)\longrightarrow
 \operatorname{Pic}(\mathrm{dP}_6).
\]
The preceding proof gives $r|_{\operatorname{Pic}^0(X)}=0$ and the local
canonical trivialization gives $r(\omega_X)=0$.  Hence, for arbitrary $L$,
\[
 \eta^*A_L\simeq r(L)^\vee.
\]
Thus the present conductor method extends from $\operatorname{Pic}^0(X)$
to every $L\in\operatorname{Pic}(X)$ exactly when $r=0$.  The still
conditional equality $H^2(X;\mathbb Z)=0$ is one sufficient route to $r=0$
by Lemma~\ref{lem:dp6-pullback}, but it is no longer needed for the
$\operatorname{Pic}^0$ statement relevant to the generic twists in CDP20.

\begin{proposition}[candidate conductor application on $\operatorname{Pic}^0$]
\label[proposition]{prop:candidate-conductor}
\Verified\ For every $L\in\operatorname{Pic}^0(X)$, putting
\[
 \mathcal F_L=T_X\otimes L,
 \qquad
 A_L=(L^\vee\otimes\omega_X)|_W,
\]
one has
\[
 H^0(W,\Omega_X^1|_W\otimes A_L)\ne0.
\]
\end{proposition}

\begin{proof}
Proposition~\ref{prop:candidate-geometry} supplies the proper smooth family,
the full reduced normal-crossings fibre and its $\mathrm{dP}_6$
normalization.  Proposition~\ref{prop:pic0-pullback} gives
$\eta^*A_L\simeq\mathcal O_{\mathrm{dP}_6}$.  The verified conductor
proposition then supplies the nonzero section.
\end{proof}

\begin{corollary}[conditional extension to all line bundles]
\label[corollary]{cor:candidate-conductor-all}
\Conditional\ If $r=0$, the conclusion of
Proposition~\ref{prop:candidate-conductor} holds for every
$L\in\operatorname{Pic}(X)$.
\end{corollary}

\subsection{The exact duality implication required}

\begin{proposition}[absolute duality on a compact Cartier fibre]
\label[proposition]{prop:cartier-fibre-duality}
\Verified\ Let $f:X\to B$ be a holomorphic map from a smooth complex
threefold to a smooth complex curve.  Suppose that
$i:W=f^{-1}(b)\hookrightarrow X$ is compact, is the full scheme-theoretic
fibre, and is an effective Cartier divisor.  For every line bundle $L$ on
$X$, put
\[
 \mathcal F_L=T_X\otimes L,
 \qquad
 A_L=(L^\vee\otimes\omega_X)|_W.
\]
Then $W$ is a pure two-dimensional Gorenstein complex space and there is a
functorial perfect pairing
\[
 H^2(W,i^*\mathcal F_L)\times
 H^0(W,\Omega_X^1|_W\otimes A_L)\longrightarrow\mathbb C.
\]
Equivalently,
\[
 H^2(W,i^*\mathcal F_L)^\vee
 \simeq H^0(W,\Omega_X^1|_W\otimes A_L).
\]
No normality hypothesis on $W$ is required.
\end{proposition}

\begin{proof}
Because $W$ is an effective Cartier divisor in the smooth threefold $X$,
each local ring of $W$ is a two-dimensional hypersurface quotient of a
regular local ring.  Thus $W$ is Cohen--Macaulay and Gorenstein.  If
$K_W^\bullet$ denotes the analytic dualizing complex, the Gorenstein
criterion and complete-intersection computation of Ramis--Ruget give
\[
 K_W^\bullet\simeq\omega_W[2]
\]
for an invertible dualizing sheaf $\omega_W$
\cite[Proposition~2 and the complete-intersection calculation,
pp.~90--91]{RamisRuget1970}.

Cartier adjunction follows by applying
$\mathcal Hom_X(-,\omega_X)$ to
\[
0\longrightarrow\mathcal O_X(-W)\longrightarrow\mathcal O_X
 \longrightarrow i_*\mathcal O_W\longrightarrow0:
\]
\[
 \omega_W=\mathcal Ext_X^1(i_*\mathcal O_W,\omega_X)
 \simeq(\omega_X\otimes\mathcal O_X(W))|_W.
\]
Since $W=f^*(b)$ as a scheme-theoretic Cartier divisor,
$\mathcal O_X(W)\simeq f^*\mathcal O_B(b)$.  Its restriction to $W$ is the
pullback of the one-dimensional vector space $\mathcal O_B(b)|_b$, so a
choice of nonzero vector gives
\[
 \mathcal O_X(W)|_W\simeq\mathcal O_W,
 \qquad \omega_W\simeq\omega_X|_W.
\]
The choice changes the final pairing by a nonzero scalar only.

Ramis--Ruget duality pairs, for coherent $\mathcal G$,
\[
 H_c^p(W,\mathcal G)
 \quad\hbox{with}\quad
 \operatorname{Ext}^{-p}_W(\mathcal G,K_W^\bullet)
\]
\cite[Theorem~2, pp.~82--83]{RamisRuget1970}.  Compactness gives
$H_c^p(W,\mathcal G)=H^p(W,\mathcal G)$.  With
$p=2$ and $\mathcal G=i^*\mathcal F_L$, the shift above gives
\[
 \operatorname{Ext}^{-2}_W(\mathcal G,K_W^\bullet)
 \simeq\operatorname{Hom}_W(\mathcal G,\omega_W).
\]
The coherent cohomology groups are finite-dimensional by Grauert finiteness
\cite{Grauert1960}; in particular their natural topologies are separated,
so the strong duality of Ramis--Ruget is an ordinary algebraic perfect
duality.  Since $\mathcal G$ is locally free,
\[
\begin{aligned}
 \operatorname{Hom}_W(\mathcal G,\omega_W)
 &=H^0(W,\mathcal G^\vee\otimes\omega_W)\\
 &\simeq
 H^0(W,\Omega_X^1|_W\otimes L^\vee|_W\otimes\omega_X|_W)\\
 &=H^0(W,\Omega_X^1|_W\otimes A_L).
\end{aligned}
\]
The pairing is cup product followed by evaluation and the trace
$H^2(W,\omega_W)\to\mathbb C$.  Precomposition on the Hom term shows the
required contravariance in $\mathcal G$.
\end{proof}

\begin{corollary}[the duality bridge used in Section~10]
\label[corollary]{cor:cartier-fibre-h2}
Under the hypotheses of the proposition,
\[
 H^0(W,\Omega_X^1|_W\otimes A_L)\ne0
 \quad\Longrightarrow\quad
 H^2(W,i^*\mathcal F_L)\ne0.
\]
\end{corollary}

\begin{corollary}[candidate fibrewise nonvanishing]
\label[corollary]{cor:candidate-fibrewise-h2}
\Verified\ For every $L\in\operatorname{Pic}^0(X)$,
\[
 H^2(W,i^*\mathcal F_L)\ne0.
\]
The duality isomorphism itself holds for every $L\in\operatorname{Pic}(X)$.
If $r=0$, the displayed nonvanishing also holds for every such $L$.
\end{corollary}

\begin{proof}
Proposition~\ref{prop:candidate-geometry} verifies that $W$ is the compact
full Cartier fibre in the smooth threefold, so the abstract duality theorem
applies.  Proposition~\ref{prop:candidate-conductor} supplies its nonzero
dual section for $L\in\operatorname{Pic}^0(X)$.  The final assertion uses
Corollary~\ref{cor:candidate-conductor-all}.
\end{proof}

\subsection{The exact top-degree base-change implication required}

\begin{proposition}[top-degree analytic base change]
\label[proposition]{prop:top-degree-base-change}
\Verified\ Let $f:X\to B$ be a proper morphism of complex spaces,
$\mathcal F$ a coherent sheaf flat over $B$ near $c\in B$, and
$X_c$ the analytic, hence full scheme-theoretic, fibre.  If
$H^3(X_c,\mathcal F|_{X_c})=0$, then the canonical map
\[
 \beta^2_c:
 (R^2f_*\mathcal F)_c\otimes_{\mathcal O_{B,c}}\mathbb C(c)
 \xrightarrow{\sim} H^2(X_c,\mathcal F|_{X_c})
\]
is an isomorphism.  In particular, a nonzero target forces
$(R^2f_*\mathcal F)_c\ne0$.
\end{proposition}

\begin{proof}
B\u{a}nic\u{a}--St\u{a}n\u{a}\c{s}il\u{a}'s base-change theorem identifies
surjectivity of the degree-$q$ fibre map with right exactness of the
corresponding direct-image functor.  Their top-degree corollary states that
if $H^q(X_c,\mathcal F|_{X_c})=0$, then these equivalent conditions hold in
degree $q-1$, and the degree-$(q-1)$ base-change maps are isomorphisms
\cite[Chapter~III, \S3, Theorem~3.4 and Corollary~3.5,
pp.~116--121]{BanicaStanasila1976}.  Taking $q=3$ gives the displayed map.
If its target is nonzero, its source is nonzero; hence the stalk itself
cannot vanish.
\end{proof}

\begin{lemma}[automatic properness and flatness over the curve]
\label[lemma]{lem:automatic-proper-flat}
\Verified\ Let $X$ be a compact connected smooth complex manifold, let $B$
be a smooth complex curve, and let $f:X\to B$ be nonconstant and
holomorphic.  Then $f$ is proper and flat.  Every locally free
$\mathcal O_X$-module is $f$-flat, and every full scheme-theoretic fibre is
an effective Cartier divisor of pure codimension one.
\end{lemma}

\begin{proof}
The underlying continuous map is proper because $X$ is compact and $B$ is
Hausdorff.  Fix $x\in X$, put $b=f(x)$, and write
$A=\mathcal O_{B,b}$ and $R=\mathcal O_{X,x}$.  The ring $A$ is a discrete
valuation ring and $R$ is a regular local analytic ring, hence an integral
domain.  The homomorphism $A\to R$ is injective: if a nonzero germ on the
curve pulled back to zero near $x$, then a neighbourhood of $x$ would map
into its discrete zero set, so $f$ would be locally constant and therefore,
by connectedness and analytic continuation, globally constant.

Thus $R$ is torsion-free as an $A$-module.  A torsion-free module over a
discrete valuation ring is flat, so $f$ is flat at $x$; since $x$ was
arbitrary, $f$ is flat.  A locally free sheaf has stalk $R^r$, which is flat
over $A$ by the same argument.  Finally, if $t$ is a local parameter at
$b$, its pullback is a nonzero element of the domain $R$ and hence a
non-zero-divisor.  Therefore the full fibre is locally cut out by one
non-zero-divisor: it is an effective Cartier divisor and is pure of
codimension one.
\end{proof}

For a full fibre of complex dimension two and a coherent sheaf, $H^3=0$.
Proposition~\ref{prop:candidate-geometry} supplies the candidate geometry,
and Lemma~\ref{lem:automatic-proper-flat} makes the locally free sheaf
$\mathcal F_L=T_X\otimes L$ flat over the base.

\begin{corollary}[candidate base change and direct-image nonvanishing]
\label[corollary]{cor:candidate-r2}
\Verified\ For every $L\in\operatorname{Pic}(X)$, the canonical map
\[
 (R^2f_*\mathcal F_L)_{p_0}\otimes\mathbb C(p_0)
 \xrightarrow{\sim}H^2(W,i^*\mathcal F_L)
\]
is an isomorphism.  Consequently, for every
$L\in\operatorname{Pic}^0(X)$,
\[
 (R^2f_*\mathcal F_L)_{p_0}\ne0,
 \qquad\text{and hence}\qquad R^2f_*(T_X\otimes L)\ne0.
\]
If $r=0$, the same nonvanishing holds for every
$L\in\operatorname{Pic}(X)$.
\end{corollary}

\begin{proof}
The candidate is proper and $\mathcal F_L$ is base-flat, while the
two-dimensional fibre has $H^3(W,i^*\mathcal F_L)=0$.  Apply
Proposition~\ref{prop:top-degree-base-change}.  The nonzero target for
$L\in\operatorname{Pic}^0(X)$ is
Corollary~\ref{cor:candidate-fibrewise-h2}; the all-line-bundle extension is
the conditional final clause of that corollary.
\end{proof}

\subsection{Present evidentiary status}

\begin{itemize}
\item The reconstructed Sections~3--6 establish the compact smooth analytic
threefold, the proper map, the full reduced normal-crossings fibre, its
$\mathrm{dP}_6$ normalization, and the local triviality of $\omega_X$ near
$W$.  These claims have now been retraced independently from the frozen
source; they use none of the open integral-topology maps.
\item For every $L\in\operatorname{Pic}^0(X)$, the pullback of $A_L$ to the
normalization is trivial.  The conductor section, compact-Cartier duality,
and top-degree base change therefore prove
$R^2f_*(T_X\otimes L)\ne0$ for every such $L$.
\item Extending this conclusion to every $L\in\operatorname{Pic}(X)$ is a
separate question: for this method it is equivalent to $r=0$.  The claimed
$H^2(X;\mathbb Z)=0$ would suffice, but still depends on the open topology
bridges.
\item The conductor section invalidates the printed cotangent-sheaf
reduction as a formal inference on this non-normal fibre.  This does not by
itself decide whether CDP20, Theorem~2.2 has another proof or whether the
candidate satisfies all of that theorem's global hypotheses.
\end{itemize}

\subsection{Remaining verification and falsification programme}

\begin{enumerate}[label=\arabic*.]
\item \textbf{All-line-bundle restriction map.} Compute
$r=\eta^*i^*$ from explicit transition functions.  A line bundle with
$r(L)\ne0$ refutes the manuscript's universal conductor argument but does
not affect the proved $\operatorname{Pic}^0$ result.  A direct proof that
$r=0$ would extend the nonvanishing to all line bundles without borrowing
the unresolved topology.

\item \textbf{Global CDP hypotheses.} Establish independently the exact
topological vanishings, $a(X)=1$, and the holomorphic-algebraic-reduction
claim used to place this $X$ under CDP20, Theorem~2.2.  The analytic
nonvanishing above is not a substitute for those global premises.

\item \textbf{Literature comparison.} Reproduce the CDP p.~692 argument
with exactly the same fibre, twist, and definition of
$\widetilde\Omega^1_S$. Determine whether an omitted hypothesis restores
the inference or whether an independent proof bypasses it. Either outcome
must be recorded without converting a proof-gap diagnosis into a theorem
verdict.
\end{enumerate}

\begin{obligation}[CDP comparison]\label[obligation]{obl:cdp-comparison}
\Open\ Produce a line-by-line diagram comparing the candidate's
ambient, conormal, and intrinsic cotangent sequences with the exact sheaves
in the cited CDP reduction. Record separately: the validity of the printed
inference, the possibility of an independent repair, and the validity of
the candidate construction.
\end{obligation}

\begin{test}[decisive outcomes]
The proved $\operatorname{Pic}^0$ chain can be refuted only by locating an
error in one of its explicit inputs: the source construction of the full
fibre, the normalization--conductor descent, the $\mathrm{dP}_6$ Picard
injection, compact-Cartier duality, or top-degree base change.  The stronger
all-line-bundle chain fails exactly when the restriction map $r$ is nonzero.
Completing proof obligation~\ref{obl:cdp-comparison} would determine whether
the printed CDP reduction admits a repair, but would still leave the global
topology obligations
\hyperref[obl:C1]{C1}--\hyperref[obl:C3]{C3},
\hyperref[obl:V1]{V1}, \hyperref[obl:MV1]{MV1}, and
\hyperref[obl:L1]{L1}.
\end{test}

\section{Confirmed local repairs kept separate from the theorem}

These corrections are mathematically definite but do not settle the central construction.

\begin{enumerate}[leftmargin=2em]
\item \Refuted\ Proposition~4.7(iv): the printed closedness justification is not valid as written.  A repair would require a complete chartwise/local-finiteness argument, which remains \Open.
\item \Refuted\ Remark~7.28 is false for a fixed triple. For example,
\[
(\ell_0,\ell_1,\ell_2)=(100,1,-1),\qquad
p=12(100)-4(1)-3(-1)=1199\ne\pm1.
\]
The intended statement must quantify existentially over $\ell_0$ for fixed admissible $(\ell_1,\ell_2)$.
\item \Verified\ Appendix Lemma~A.3 prints cohomological arrows while labelling them as homology. The sequence must be written
\[
H^k(W)\longrightarrow H^k(\widetilde W)\oplus H^k(D)
\longrightarrow H^k(\widetilde D)\longrightarrow H^{k+1}(W),
\]
followed by universal coefficients when homology ranks are required.
\item \Verified\ Appendix Remark~A.6 must read $\pi_1(N_0)=\bar\Lambda=\pi_1(W)$, not $\pi_1(N_0)=\Lambda$.
\item \Verified\ The modular assembler has a page-71 portability defect under the tested MiK\TeX{} route; the canonical monolithic source renders correctly. This is a build defect, not a mathematical counterexample.
\end{enumerate}

\section{Research log}

\subsection*{Round 0 --- private workbench initialization}

\begin{itemize}[leftmargin=1.8em]
\item Preserved the completed 121-page annotated audit as \Source{audit_baseline/annotated_audit.tex} and the 108-page reconstruction as \Source{audit_baseline/clean_transcription.tex}.
\item Converted the review objections into named mathematical obligations with explicit maps, bases, signs, and falsification tests.
\item Proved the exact conditional endpoint: once the geometric maps are identified with the printed matrices, the verified algebra gives the standard smooth $6$-sphere.
\item Kept the conductor mechanism separate from both the existence of $X$ and the truth or falsity of the published CDP theorem.
\end{itemize}

\subsection*{Next calculations}

\begin{enumerate}[leftmargin=1.8em]
\item Extract the actual coordinate formulas for the three collar inclusions and build a complete orientation/sign table.
\item Build a finite integral chain model for the cusp central fibre $W$ and compare its specialization map with Appendix~B in primitive bases.
\item Compute the Leray $d_2$ geometrically and compare it to $d_2(12\gamma)=\pm p\omega$ without importing the Mayer--Vietoris normalization.
\item Compute the restriction map
$r=\eta^*i^*:\operatorname{Pic}(X)\to\operatorname{Pic}(\mathrm{dP}_6)$
from transition functions, and complete the line-by-line CDP cotangent-sheaf
comparison without using the unresolved topology.
\end{enumerate}

\subsection*{Round 1 --- 26 August 2026: typed workbench release}

\begin{itemize}[leftmargin=1.8em]
\item Rebuilt the workbench as a new clean tree with the exact title
``S6 Proof Repair and Disproof Workbench --- Private''.  The public audit
project is not an input to the upload workflow and remains untouched.
\item Kept the terminal algebra \Verified{}, but kept the global topology
\Conditional{} until C1--C3, V1, MV1, and L1 identify the printed matrices with
actual geometric maps.  No theorem-level CDP verdict is inferred from a
proof-step defect alone.
\item Added a reader-facing literature audit with DOI-resolvable references
for CDP20/CDP98, analytic duality and base change, nearby cycles, topology,
and the Leiden Declaration.  Local SGA TeX was used as a comparison layer;
third-party files are not bundled.
\item Added the bounded Fable/Jacobian provenance question to the
comparative workstream.  Any source path not independently located remains
\Open{} rather than being presented as a recovered transcript.
\end{itemize}

\subsection*{Round 2 --- 26 August 2026: duality and base-change bridges}

\begin{itemize}[leftmargin=1.8em]
\item Proved the abstract compact-Cartier-fibre duality statement from
Ramis--Ruget's analytic dualizing complex and trace pairing, Cartier
adjunction, and coherent-cohomology finiteness.  Normality of the fibre is
not required.
\item Promoted the already complete normalization--conductor calculation to
\Verified{} and separated it from its \Conditional{} application to the
candidate.  Also isolated the elementary $\mathrm{dP}_6$ Picard pullback
lemma: $H^2(X;\mathbb Z)=0$ forces every line bundle from $X$ to pull back
trivially to the normalization.
\item At this round, promoted only the abstract theorem to \Verified{} and
left its candidate application \Conditional{} on smoothness, compactness,
the full scheme-theoretic fibre, and Cartierness.  Round~3 below supersedes
that interim status after independently retracing Sections~3--6.
\item Proved the required abstract top-degree base-change isomorphism from
the degree-$3$ vanishing corollary of the analytic base-change theorem.
Also proved that properness, flatness, Cartierness of the full fibre, and
$f$-flatness of a vector bundle are automatic for a nonconstant map from a
compact connected smooth complex manifold to a smooth curve.
\item At this round, kept the candidate instantiation \Open{} pending a chart
audit.  Round~3 below records the completed source reconstruction and status
promotion.
\item Reconciled the conditional topology endpoint with the six exact bridge
obligations C1--C3, V1, MV1, and L1; no four-item shorthand now substitutes
for that dependency set.
\end{itemize}

\subsection*{Round 3 --- 26 August 2026: candidate geometry and the
$\operatorname{Pic}^0$ strengthening}

\begin{itemize}[leftmargin=1.8em]
\item Independently retraced Theorem~3.4, Theorem~4.5,
Proposition~4.6, Proposition~4.7(i)--(iii), Theorem~5.4, and
Theorem~6.2 in the frozen transcription.  This closes the analytic action,
quotient, global gluing, full-fibre, $\mathrm{dP}_6$ normalization, and local
canonical-triviality gates.  It does not close the cusp-collapse or integral
topology maps.
\item Proved a topology-free Picard lemma: for every
$L\in\operatorname{Pic}^0(X)$, the line bundle
$A_L=(L^\vee\otimes\omega_X)|_W$ pulls back trivially to
$\mathrm{dP}_6$.  The proof uses $c_1(L)=0$, injectivity of the Chern-class
map on $\mathrm{dP}_6$, and the nowhere-vanishing three-form on $N_0$.
\item Combined that lemma with the verified conductor section, analytic
duality, and top-degree base change to obtain
$R^2f_*(T_X\otimes L)\ne0$ for every
$L\in\operatorname{Pic}^0(X)$.  This conclusion no longer borrows
$H^2(X;\mathbb Z)=0$.
\item Isolated the exact remaining all-line-bundle question as
$r=\eta^*i^*=0$.  The open global topology would imply it, but the
$\operatorname{Pic}^0$ conclusion needed for generic CDP twists is already
independent of that route.
\end{itemize}

\section{Literature audit and source gate}

This is a publication-facing literature record, not a list of private chat
notes.  The CDP comparison is against Campana--Demailly--Peternell's 2020
paper \cite[Theorem~2.2 and Corollary~2.3, p.~680]{CDP2020}, with the 1998
paper \cite[Lemma~3.2, pp.~84--85]{CDP1998} retained because the later argument imports its
semidirect-product calculation.  The exact theorem and page locators are
given in the bibliography below; the DOI links are the reader-facing source
gate.

For the analytic part, the workbench uses the absolute duality results of
Ramis--Ruget \cite[Theorems~1--2 and Proposition~2, pp.~81--91]{RamisRuget1970}
and the relative analytic reference of
Ramis--Ruget--Verdier \cite{RamisRugetVerdier1971}, together with
Bănică--Stănășilă's cohomology-and-base-
change criteria \cite{BanicaStanasila1976}.  The latter is cited at Chapter~III, §3, Theorems~3.4 and
Corollaries~3.5 and~3.10, and §4, Theorem~4.12.  The Stacks Project is used
only as an algebraic cross-check, not as a substitute for the analytic
theorem.  Nearby-cycle formalism is compared with SGA~7~I/II and with the
analytic treatment cited in the bibliography; none of those general
references identifies the candidate's printed matrices without a
candidate-specific map calculation.

The local literature corpus was also checked before writing this section.
The relevant SGA~6 and SGA~7 TeX paths are recorded verbatim in
\Source{research/README.md} and
\Source{comparative_context/README.md}; they were read as a local comparison
layer.  No third-party source files or PDFs are copied into this project.  The
public citations remain the canonical way for a reader to check the literature.

The mathematical role of each source is deliberately narrow:

\begin{itemize}[leftmargin=1.8em]
\item CDP20 supplies the theorem-level target and the published proof steps
under audit; CDP98 supplies the earlier split-extension calculation.
\item Ramis--Ruget, Ramis--Ruget--Verdier, and Bănică--Stănășilă supply the
analytic duality and base-change statements whose hypotheses are still
checked explicitly in Section~10.
\item SGA~7 and the analytic nearby-cycle references supply formal
frameworks \cite{SGA7I1972,SGA7II1973}, not the candidate-specific integral identifications.
\item Hatcher, Smale, and Kervaire--Milnor supply the final recognition
theorems \cite{Hatcher2002,Smale1961,KervaireMilnor1963} only after the geometric maps have been certified.
\item The Leiden Declaration is recorded as a methodological disclosure
reference \cite{Leiden2026}; it is not mathematical evidence for any claim.
\end{itemize}

No literature result is silently promoted from a general theorem to a fact
about the candidate.  Any such promotion must appear as a typed claim in the
ledgers and carry its own map, hypothesis, and falsification test.

\begin{thebibliography}{99}

\bibitem{Leiden2026}
J.~Alper, M.~Barany, A.~Chavarri Villarello, S.~Dahmen, W.~Dean,
K.~Ganapathy, M.~Harris, D.~Holmes, M.~Jamnik, S.~Kelk, B.~Kra,
U.~Martin, B.~Naskr\k{e}cki, R.~Ochigame, J.~Portegies, and J.~Schmitt,
\emph{Leiden Declaration on Artificial Intelligence and Mathematics},
2026,
\href{https://doi.org/10.5281/zenodo.20302944}{doi:10.5281/zenodo.20302944}.

\bibitem{CDP2020}
F.~Campana, J.-P.~Demailly, and T.~Peternell,
\emph{The algebraic dimension of compact complex threefolds with vanishing
second Betti number},
Compos. Math. \textbf{156} (2020), 679--696,
\href{https://doi.org/10.1112/S0010437X19007802}{doi:10.1112/S0010437X19007802}.

\bibitem{CDP1998}
F.~Campana, J.-P.~Demailly, and T.~Peternell,
\emph{The algebraic dimension of compact complex threefolds with vanishing
second Betti number},
Compos. Math. \textbf{112} (1998), 77--91,
\href{https://doi.org/10.1023/A:1000313214795}{doi:10.1023/A:1000313214795};
\href{https://arxiv.org/abs/math/9607215}{arXiv:math/9607215}.

\bibitem{RamisRuget1970}
J.-P.~Ramis and G.~Ruget,
\emph{Complexe dualisant et th\'eor\`emes de dualit\'e en g\'eom\'etrie
analytique complexe},
Publ. Math. IH\'ES \textbf{38} (1970), 77--91,
\href{https://doi.org/10.1007/BF02684652}{doi:10.1007/BF02684652}.

\bibitem{RamisRugetVerdier1971}
J.-P.~Ramis, G.~Ruget, and J.-L.~Verdier,
\emph{Dualit\'e relative en g\'eom\'etrie analytique complexe},
Invent. Math. \textbf{13} (1971), 261--283,
\href{https://doi.org/10.1007/BF01406078}{doi:10.1007/BF01406078};
\href{https://eudml.org/doc/142096}{EuDML record}.

\bibitem{BanicaStanasila1976}
C.~B\u{a}nic\u{a} and O.~St\u{a}n\u{a}\c{s}il\u{a},
\emph{Algebraic Methods in the Global Theory of Complex Spaces},
Editura Academiei and John Wiley \& Sons, Bucharest and London, 1976;
  Chapter~III, \S3, Theorem~3.4 and Corollaries~3.5 and~3.10, pp.~116--123;
  Chapter~III, \S4, Theorem~4.12, pp.~134--135,
  \href{https://djvu.online/file/2hQiOnz8WuoQk}{publicly accessible scan}.

\bibitem{Grauert1960}
H.~Grauert, \emph{Ein Theorem der analytischen Garbentheorie und die
Modulr\"aume komplexer Strukturen}, Publ. Math. IH\'ES \textbf{5} (1960),
5--64, \href{https://doi.org/10.1007/BF02684746}{doi:10.1007/BF02684746}.

\bibitem{Lee2024}
J.~M.~Lee,
\emph{Introduction to Complex Manifolds},
Graduate Studies in Mathematics \textbf{244}, American Mathematical Society,
Providence, RI, 2024; Theorem~9.37,
\href{https://bookstore.ams.org/gsm-244}{publisher page}.

\bibitem{Hatcher2002}
A.~Hatcher,
\emph{Algebraic Topology},
Cambridge University Press, Cambridge, 2002; Theorems~4.5 and~4.32 and
Corollary~4.33,
\href{https://pi.math.cornell.edu/~hatcher/AT/ATpage.html}{author's electronic edition}.

\bibitem{Bredon1997}
G.~E.~Bredon,
\emph{Sheaf Theory}, second edition,
Graduate Texts in Mathematics \textbf{170}, Springer, New York, 1997;
Chapter~VI,
\href{https://doi.org/10.1007/978-1-4612-0647-7}{doi:10.1007/978-1-4612-0647-7}.

\bibitem{Smale1961}
S.~Smale,
\emph{Generalized Poincar\'e's conjecture in dimensions greater than four},
Ann. of Math. (2) \textbf{74} (1961), 391--406,
\href{https://doi.org/10.2307/1970239}{doi:10.2307/1970239}.

\bibitem{KervaireMilnor1963}
M.~A.~Kervaire and J.~W.~Milnor,
\emph{Groups of homotopy spheres: I},
Ann. of Math. (2) \textbf{77} (1963), 504--537,
\href{https://doi.org/10.2307/1970128}{doi:10.2307/1970128}.

\bibitem{Hartshorne1977}
R.~Hartshorne,
\emph{Algebraic Geometry},
Graduate Texts in Mathematics \textbf{52}, Springer, New York, 1977.

\bibitem{BirkenhakeLange1999}
C.~Birkenhake and H.~Lange,
\emph{Complex Tori},
Progress in Mathematics \textbf{177}, Birkh\"auser, Boston, 1999;
Chapter~2, \S6,
\href{https://doi.org/10.1007/978-1-4612-1566-0}{doi:10.1007/978-1-4612-1566-0}.

\bibitem{PareschiPopaSchnell2017}
G.~Pareschi, M.~Popa, and C.~Schnell,
\emph{Hodge modules on complex tori and generic vanishing for compact
K\"ahler manifolds},
Geom. Topol. \textbf{21} (2017), 2419--2460; \S8, p.~2437,
\href{https://doi.org/10.2140/gt.2017.21.2419}{doi:10.2140/gt.2017.21.2419}.

\bibitem{Fulton1993}
W.~Fulton,
\emph{Introduction to Toric Varieties},
Annals of Mathematics Studies \textbf{131}, Princeton University Press,
Princeton, 1993.

\bibitem{CoxLittleSchenck2011}
D.~A.~Cox, J.~B.~Little, and H.~K.~Schenck,
\emph{Toric Varieties},
Graduate Studies in Mathematics \textbf{124}, American Mathematical Society,
Providence, RI, 2011; Theorem~12.1.10,
\href{https://doi.org/10.1090/gsm/124}{doi:10.1090/gsm/124}.

\bibitem{Danilov1978}
V.~I.~Danilov,
\emph{The geometry of toric varieties},
Russian Math. Surveys \textbf{33} (1978), no.~2, 97--154; \S9,
\href{https://doi.org/10.1070/RM1978v033n02ABEH002305}{doi:10.1070/RM1978v033n02ABEH002305}.

\bibitem{Serre1973}
J.-P.~Serre,
\emph{A Course in Arithmetic},
Graduate Texts in Mathematics \textbf{7}, Springer, New York, 1973;
Chapter~VII,
\href{https://doi.org/10.1007/978-1-4684-9884-4}{doi:10.1007/978-1-4684-9884-4}.

\bibitem{Brown1982}
K.~S.~Brown,
\emph{Cohomology of Groups},
Graduate Texts in Mathematics \textbf{87}, Springer, New York, 1982.

\bibitem{SGA7I1972}
A.~Grothendieck, M.~Raynaud, and D.~S.~Rim (eds.),
\emph{Groupes de monodromie en g\'eom\'etrie alg\'ebrique. I},
S\'eminaire de G\'eom\'etrie Alg\'ebrique du Bois-Marie 1967--1969,
Lecture Notes in Mathematics \textbf{288}, Springer, Berlin, 1972,
\href{https://doi.org/10.1007/BFb0068688}{doi:10.1007/BFb0068688}.

\bibitem{SGA7II1973}
P.~Deligne and N.~Katz (eds.),
\emph{Groupes de monodromie en g\'eom\'etrie alg\'ebrique. II},
S\'eminaire de G\'eom\'etrie Alg\'ebrique du Bois-Marie 1967--1969,
Lecture Notes in Mathematics \textbf{340}, Springer, Berlin, 1973.

\bibitem{Dimca2004}
A.~Dimca,
\emph{Sheaves in Topology},
Universitext, Springer, Berlin, 2004; Definition~4.2.1,
Propositions~4.2.2 and~4.2.11, Theorem~2.3.26, and Corollary~2.3.4,
\href{https://doi.org/10.1007/978-3-642-18868-8}{doi:10.1007/978-3-642-18868-8}.

\bibitem{KashiwaraSchapira1990}
M.~Kashiwara and P.~Schapira,
\emph{Sheaves on Manifolds},
Grundlehren der mathematischen Wissenschaften \textbf{292}, Springer,
Berlin, 1990; \S8.6,
\href{https://doi.org/10.1007/978-3-662-02661-8}{doi:10.1007/978-3-662-02661-8}.

\bibitem{SGAEnglish2026}
A.~Grothendieck and the SGA contributors,
\emph{SGA in English: maintained working edition}, 2026,
\href{https://doi.org/10.5281/zenodo.21925497}{doi:10.5281/zenodo.21925497};
\href{https://github.com/KokunoYumeto/sga-en}{source repository}.
This working edition is cited only as a searchable reading layer; the
original volumes remain the authority.

\bibitem{StacksProject}
The Stacks Project Authors,
\emph{The Stacks Project},
\href{https://stacks.math.columbia.edu}{stacks.math.columbia.edu},
especially Tags
\href{https://stacks.math.columbia.edu/tag/0AUW}{0AUW} and
\href{https://stacks.math.columbia.edu/tag/0EX2}{0EX2}.

\end{thebibliography}


\end{document}
